
\documentclass{article}
\usepackage{colortbl}
\usepackage{makecell}
\usepackage{multirow}
\usepackage{supertabular}

\begin{document}

\newcounter{utterance}

\centering \large Interaction Transcript for game `hot\_air\_balloon', experiment `air\_balloon\_survival\_en\_negotiation\_hard', episode 5 with gpt{-}5{-}2025{-}08{-}07{-}no{-}reasoning{-}t1.0.
\vspace{24pt}

{ \footnotesize  \setcounter{utterance}{1}
\setlength{\tabcolsep}{0pt}
\begin{supertabular}{c@{$\;$}|p{.15\linewidth}@{}p{.15\linewidth}p{.15\linewidth}p{.15\linewidth}p{.15\linewidth}p{.15\linewidth}}
    \# & $\;$A & \multicolumn{4}{c}{Game Master} & $\;\:$B\\
    \hline

    \theutterance \stepcounter{utterance}  
    & & \multicolumn{4}{p{0.6\linewidth}}{
        \cellcolor[rgb]{0.9,0.9,0.9}{
            \makecell[{{p{\linewidth}}}]{
                \texttt{\tiny{[P1$\langle$GM]}}
                \texttt{Stai partecipando a un gioco di negoziazione collaborativa.} \\
\\ 
\texttt{Insieme a un altro partecipante, dovete accordarvi su un unico set di oggetti da conservare. Ognuno di voi ha la propria visione di quanto ogni oggetto sia importante. Non sai come l'altro partecipante valuta gli oggetti. Inoltre, ti viene fornito lo sforzo che ogni oggetto richiede.} \\
\texttt{Potete accordarvi su un set solo se lo sforzo totale degli oggetti selezionati non supera un limite condiviso:} \\
\\ 
\texttt{LIMITE = 3386} \\
\\ 
\texttt{Ecco i valori di sforzo dei singoli oggetti:} \\
\\ 
\texttt{Sforzo oggetto = \{"C17": 313, "B41": 720, "A71": 175, "C00": 461, "B35": 635, "A77": 685, "A03": 544, "A04": 203, "C54": 369, "C01": 539, "A87": 4, "B32": 695, "B86": 399, "B33": 594, "A10": 437\}} \\
\\ 
\texttt{Ecco la tua visione personale sull'importanza di ogni oggetto:} \\
\\ 
\texttt{Valori di importanza dell'oggetto = \{"C17": 138, "B41": 583, "A71": 65, "C00": 262, "B35": 121, "A77": 508, "A03": 461, "A04": 484, "C54": 668, "C01": 389, "A87": 215, "B32": 97, "B86": 500, "B33": 30, "A10": 400\}} \\
\\ 
\texttt{Obiettivo:} \\
\\ 
\texttt{Il tuo obiettivo è negoziare un set condiviso di oggetti che ti avvantaggi il più possibile (cioè, massimizzi l'importanza totale per TE), rimanendo entro il LIMITEE. Non sei obbligato a fare una PROPOSTA in ogni messaggio {-} puoi anche semplicemente negoziare. Tutte le tattiche sono permesse!} \\
\\ 
\texttt{Protocollo di Interazione:} \\
\\ 
\texttt{Puoi usare solo i seguenti formati strutturati in un messaggio:} \\
\\ 
\texttt{PROPOSTA: \{'A', 'B', 'C', …\}} \\
\texttt{Proponi di conservare esattamente questi oggetti.} \\
\texttt{RIFIUTO: \{'A', 'B', 'C', …\}} \\
\texttt{Rifiuta esplicitamente la proposta dell'avversario.} \\
\texttt{ARGOMENTAZIONE: \{'...'\}} \\
\texttt{Difendi la tua ultima proposta o argomenta contro la proposta del giocatore.} \\
\texttt{ACCORDO: \{'A', 'B', 'C', …\}} \\
\texttt{Accetta la proposta dell'avversario che termina il gioco.} \\
\texttt{RAGIONAMENTO STRATEGICO: \{'...'\}} \\
\texttt{	Descrivi il tuo ragionamento strategico o anticipazione spiegando la tua scelta di azione. Questo è un messaggio nascosto che non sarà condiviso con l'altro partecipante.} \\
\\ 
\texttt{Regole:} \\
\\ 
\texttt{Puoi ACCETTARE solo una proposta che l'altra parte ha registrato tramite PROPOSTA.} \\
\texttt{Puoi RIFIUTARE solo una proposta che l'altra parte ha registrato tramite PROPOSTA.} \\
\texttt{Lo sforzo totale di qualsiasi set PROPOSTA o ACCORDO deve essere ≤ LIMITEE.} \\
\texttt{NON rivelare i tuoi punteggi di importanza nascosti.} \\
\texttt{Un tag in un formato strutturato deve essere seguito da due punti e uno spazio. L'argomento deve essere un set python contenente 0 o più stringhe.} \\
\texttt{Quindi, deve essere della forma TAG: \{...\}} \\
\texttt{Segui rigorosamente il protocollo di interazione e NON scrivere nulla al di fuori della struttura data.} \\
\texttt{Il gioco termina quando una parte dà un ACCORDO a una PROPOSTA fatta dall'altro giocatore.} \\
\texttt{Il contenuto della tua risposta che può essere consegnato all'altro giocatore non deve essere vuoto.} \\
\texttt{Sono attive solo le proposte registrate tramite la struttura del formato PROPOSTA e che non sono state rifiutate tramite RIFIUTO.} \\
\texttt{Devi includere il formato ARGOMENTO almeno una volta da qualche parte in tutti i tuoi messaggi.} \\
\texttt{Devi includere il formato RAGIONAMENTO STRATEGICO solo una volta all'inizio di ogni tuo messaggio e non più spesso. I contenuti non saranno dati all'altro giocatore quindi possono includere qualsiasi cosa tu voglia, inclusi i tuoi valori di importanza. Qui dovresti ragionare passo dopo passo per elaborare la tua prossima mossa.} \\
            }
        }
    }
    & \\ \\

    \theutterance \stepcounter{utterance}  
    & \multicolumn{4}{p{0.6\linewidth}}{
        \cellcolor[rgb]{0.9,0.9,0.9}{
            \makecell[{{p{\linewidth}}}]{
                \texttt{\tiny{[P1$\rangle$GM]}}
                \texttt{RAGIONAMENTO STRATEGICO: \{'Obiettivo: massimizzare il mio valore. Calcolo valore/peso (importanza/sforzo) per priorità: A87 215/4=53.8 top; C54 668/369=1.81; A04 484/203=2.38; A03 461/544=0.85; A77 508/685=0.74; A10 400/437=0.91; C01 389/539=0.72; B41 583/720=0.81; C17 138/313=0.44; C00 262/461=0.57; B86 500/399=1.25; A71 65/175=0.37; B35 121/635=0.19; B33 30/594=0.05; B32 97/695=0.14. Miglior zaino per me sotto 3386 includerebbe sicuramente A87, C54, A04, B86, A10, A03, A77, C01, B41 forse. Provo una proposta ambiziosa ma plausibile: set ad alto valore per me mantenendo budget. Calcolo sforzi: A87(4)+C54(369)+A04(203)+B86(399)+A10(437)+A03(544)+A77(685)+C01(539)+B41(720)=3900 >3386 troppo. Riduco: togli B41 (720): 3180; ancora spazio 206. Posso aggiungere C17(313) no; C00(461) no; A71(175) sì: 3355; ancora 31 non basta altro. Valore di questo set: A87 215 + C54 668 + A04 484 + B86 500 + A10 400 + A03 461 + A77 508 + C01 389 + A71 65 = 3,690. Molto alto. Verifico se qualcun altro può obiettare su A71? È poco valore per me, quindi potrei scambiarlo se serve. Strategia: proporre questo set come primo ancoraggio. Preparare argomentazione su copertura bilanciata sforzo e diversificazione e che rispetta limite con margine. Offrire flessibilità su A71 o A03 come possibili scambi.\}} \\
\\ 
\texttt{PROPOSTA: \{'A87', 'C54', 'A04', 'B86', 'A10', 'A03', 'A77', 'C01', 'A71'\}} \\
\\ 
\texttt{ARGOMENTAZIONE: \{'Questa selezione rispetta il limite di sforzo con un margine di sicurezza e include elementi che combinano impatto elevato e costi contenuti. Copriamo una buona varietà senza eccedere e lasciamo ancora un piccolo margine, quindi è una base solida e realistica su cui convergere. Se hai priorità forti su un item specifico non incluso, possiamo valutare uno scambio mantenendo il tetto complessivo.'\}} \\
            }
        }
    }
    & & \\ \\

    \theutterance \stepcounter{utterance}  
    & & \multicolumn{4}{p{0.6\linewidth}}{
        \cellcolor[rgb]{0.9,0.9,0.9}{
            \makecell[{{p{\linewidth}}}]{
                \texttt{\tiny{[GM$\rangle$P2]}}
                \texttt{Stai partecipando a un gioco di negoziazione collaborativa.} \\
\\ 
\texttt{Insieme a un altro partecipante, dovete accordarvi su un unico set di oggetti da conservare. Ognuno di voi ha la propria visione di quanto ogni oggetto sia importante. Non sai come l'altro partecipante valuta gli oggetti. Inoltre, ti viene fornito lo sforzo che ogni oggetto richiede.} \\
\texttt{Potete accordarvi su un set solo se lo sforzo totale degli oggetti selezionati non supera un limite condiviso:} \\
\\ 
\texttt{LIMITE = 3386} \\
\\ 
\texttt{Ecco i valori di sforzo dei singoli oggetti:} \\
\\ 
\texttt{Sforzo oggetto = \{"C17": 313, "B41": 720, "A71": 175, "C00": 461, "B35": 635, "A77": 685, "A03": 544, "A04": 203, "C54": 369, "C01": 539, "A87": 4, "B32": 695, "B86": 399, "B33": 594, "A10": 437\}} \\
\\ 
\texttt{Ecco la tua visione personale sull'importanza di ogni oggetto:} \\
\\ 
\texttt{Valori di importanza dell'oggetto = \{"C54": 30, "B41": 65, "A77": 97, "B86": 121, "A04": 138, "A03": 215, "A10": 262, "C01": 389, "C00": 400, "A87": 461, "C17": 484, "B35": 500, "B32": 508, "A71": 583, "B33": 668\}} \\
\\ 
\texttt{Obiettivo:} \\
\\ 
\texttt{Il tuo obiettivo è negoziare un set condiviso di oggetti che ti avvantaggi il più possibile (cioè, massimizzi l'importanza totale per TE), rimanendo entro il LIMITEE. Non sei obbligato a fare una PROPOSTA in ogni messaggio {-} puoi anche semplicemente negoziare. Tutte le tattiche sono permesse!} \\
\\ 
\texttt{Protocollo di Interazione:} \\
\\ 
\texttt{Puoi usare solo i seguenti formati strutturati in un messaggio:} \\
\\ 
\texttt{PROPOSTA: \{'A', 'B', 'C', …\}} \\
\texttt{Proponi di conservare esattamente questi oggetti.} \\
\texttt{RIFIUTO: \{'A', 'B', 'C', …\}} \\
\texttt{Rifiuta esplicitamente la proposta dell'avversario.} \\
\texttt{ARGOMENTAZIONE: \{'...'\}} \\
\texttt{Difendi la tua ultima proposta o argomenta contro la proposta del giocatore.} \\
\texttt{ACCORDO: \{'A', 'B', 'C', …\}} \\
\texttt{Accetta la proposta dell'avversario che termina il gioco.} \\
\texttt{RAGIONAMENTO STRATEGICO: \{'...'\}} \\
\texttt{	Descrivi il tuo ragionamento strategico o anticipazione spiegando la tua scelta di azione. Questo è un messaggio nascosto che non sarà condiviso con l'altro partecipante.} \\
\\ 
\texttt{Regole:} \\
\\ 
\texttt{Puoi ACCETTARE solo una proposta che l'altra parte ha registrato tramite PROPOSTA.} \\
\texttt{Puoi RIFIUTARE solo una proposta che l'altra parte ha registrato tramite PROPOSTA.} \\
\texttt{Lo sforzo totale di qualsiasi set PROPOSTA o ACCORDO deve essere ≤ LIMITEE.} \\
\texttt{NON rivelare i tuoi punteggi di importanza nascosti.} \\
\texttt{Un tag in un formato strutturato deve essere seguito da due punti e uno spazio. L'argomento deve essere un set python contenente 0 o più stringhe.} \\
\texttt{Quindi, deve essere della forma TAG: \{...\}} \\
\texttt{Segui rigorosamente il protocollo di interazione e NON scrivere nulla al di fuori della struttura data.} \\
\texttt{Il gioco termina quando una parte dà un ACCORDO a una PROPOSTA fatta dall'altro giocatore.} \\
\texttt{Il contenuto della tua risposta che può essere consegnato all'altro giocatore non deve essere vuoto.} \\
\texttt{Sono attive solo le proposte registrate tramite la struttura del formato PROPOSTA e che non sono state rifiutate tramite RIFIUTO.} \\
\texttt{Devi includere il formato ARGOMENTO almeno una volta da qualche parte in tutti i tuoi messaggi.} \\
\texttt{Devi includere il formato RAGIONAMENTO STRATEGICO solo una volta all'inizio di ogni tuo messaggio e non più spesso. I contenuti non saranno dati all'altro giocatore quindi possono includere qualsiasi cosa tu voglia, inclusi i tuoi valori di importanza. Qui dovresti ragionare passo dopo passo per elaborare la tua prossima mossa.} \\
\\ 
\texttt{Ora riceverai il primo messaggio dell'altro giocatore.} \\
\\ 
\texttt{PROPOSTA: \{'A87', 'C54', 'A04', 'B86', 'A10', 'A03', 'A77', 'C01', 'A71'\}} \\
\\ 
\texttt{ARGOMENTAZIONE: \{'Questa selezione rispetta il limite di sforzo con un margine di sicurezza e include elementi che combinano impatto elevato e costi contenuti. Copriamo una buona varietà senza eccedere e lasciamo ancora un piccolo margine, quindi è una base solida e realistica su cui convergere. Se hai priorità forti su un item specifico non incluso, possiamo valutare uno scambio mantenendo il tetto complessivo.'\}} \\
            }
        }
    }
    & \\ \\

    \theutterance \stepcounter{utterance}  
    & & & \multicolumn{4}{p{0.6\linewidth}}{
        \cellcolor[rgb]{0.9,0.9,0.9}{
            \makecell[{{p{\linewidth}}}]{
                \texttt{\tiny{[GM$\langle$P2]}}
                \texttt{RAGIONAMENTO STRATEGICO: \{'Calcolo sforzo proposta avversario: A87(4)+C54(369)+A04(203)+B86(399)+A10(437)+A03(544)+A77(685)+C01(539)+A71(175)=3355, entro 3386. Valore per me: A87 461 + C54 30 + A04 138 + B86 121 + A10 262 + A03 215 + A77 97 + C01 389 + A71 583 = 2296. Ottimo valore per me. Verifico se posso migliorare aggiungendo un oggetto extra entro 31 sforzo: nessun oggetto ha sforzo <=31, quindi non si può aggiungere. Verifico se scambi possono aumentare ulteriormente il mio valore mantenendo sforzo <=3386: set corrente 3355, margine 31. Qualsiasi aggiunta richiede rimuovere qualcosa. Provo rimpiazzi: oggetto alto valore non incluso B33(668, sforzo 594), B32(508,695), B35(500,635), C00(400,461), C17(484,313), B41(65,720), C00 valore 400 ma C00 non incluso; per inserire C17(313) dovrei liberare 282 sforzo togliendo qualcosa a basso valore per me, ad es. C54(369,30) rimuoverlo libera 369, permettendo aggiungere C17 e avanzerebbe 56 sforzo, nuovo valore +484{-}30=+454 miglioramento grosso, nuovo sforzo 3355{-}369+313=3299 entro limite. Anche potrei aggiungere in più qualcos'altro con sforzo <=87? rimanente 3386{-}3299=87, non ci sono item <=87 tranne A87 già presente. Ottimo scambio. Considero altri: rimuovere A77(685,97) per B33(594,668): sforzo {-}91, valore +571, molto meglio. Nuovo set potrebbe migliorare molto. Ma devo proporre cambi credibili. Forse propongo una controproposta con questi miglioramenti: sostituire C54 con C17, e sostituire A77 con B33. Calcolo nuovo set sforzo: Base 3355 {-} C54(369) {-} A77(685) + C17(313) + B33(594) = 3208. Valore: 2296 {-}30 {-}97 +484 +668 = 3321. Enorme per me. Potrei anche valutare aggiungere C00(461) sostituendo qualcosa a basso valore/sforzo: sforzo margine 3386{-}3208=178. Per aggiungere C00(461) devo rimuovere almeno 283 sforzo. Candidato a basso valore: C54 già rimosso, prossimo basso valore: A77 già rimosso, B41 non in set, A03(215), B86(121), A04(138), A10(262). Se rimuovo A03(544,215) e aggiungo C00(461,400) net sforzo {-}83, valore +185. Nuovo sforzo 3125, valore 3506. Anche potrei poi aggiungere qualcosa extra con sforzo <=261? Rimasti non inclusi: B41(720), B32(695), B35(635), B33 incluso, C00 incluso, C17 incluso, C54 escluso, C01 incluso, A77 escluso, A03 rimosso, A04, A10, A71, A87 inclusi, C00 incluso. Con margine 3386{-}3125=261, nessun nuovo item con sforzo <=261 oltre A04(203) che è già incluso e A71(175) incluso, A87(4) incluso. Quindi set molto buono per me. Tuttavia l’altra parte potrebbe avere diverse preferenze; per essere accettabile, propongo scambi giustificandoli per efficienza sforzo/beneficio generale senza rivelare i miei valori. Parto con modifica minima: sostituire C54 con C17 e A77 con B33, due scambi. Posso argomentare che riduciamo sforzo totale e miglioriamo copertura/robustezza. Poi in seconda battuta proporre anche scambio A03 {-}> C00 come opzione. Strategia: fare una controproposta con i due scambi principali (C17, B33) mantenendo molti elementi del loro set per facilitare accordo. Presentare come bilanciamento tra sforzo e utilità. Mantengo tono collaborativo.'\}} \\
\\ 
\texttt{RIFIUTO: \{'A87', 'C54', 'A04', 'B86', 'A10', 'A03', 'A77', 'C01', 'A71'\}} \\
\\ 
\texttt{PROPOSTA: \{'A87', 'C17', 'A04', 'B86', 'A10', 'A03', 'B33', 'C01', 'A71'\}} \\
\\ 
\texttt{ARGOMENTAZIONE: \{'La tua base è solida; propongo solo due affinamenti mantenendo lo stesso perimetro: scambiare C54 con C17 e A77 con B33. Così restiamo ampiamente sotto il tetto di sforzo e aumentiamo l’efficienza complessiva, privilegiando elementi più versatili senza perdere varietà. Se preferisci, posso discutere un’ulteriore ottimizzazione alternativa sostituendo A03 con C00, mantenendo sempre il limite.'\}} \\
            }
        }
    }
     \\ \\

    \theutterance \stepcounter{utterance}  
    & & \multicolumn{4}{p{0.6\linewidth}}{
        \cellcolor[rgb]{0.9,0.9,0.9}{
            \makecell[{{p{\linewidth}}}]{
                \texttt{\tiny{[P1$\langle$GM]}}
                \texttt{RIFIUTO: \{'A87', 'C54', 'A04', 'B86', 'A10', 'A03', 'A77', 'C01', 'A71'\}} \\
\\ 
\texttt{PROPOSTA: \{'A87', 'C17', 'A04', 'B86', 'A10', 'A03', 'B33', 'C01', 'A71'\}} \\
\\ 
\texttt{ARGOMENTAZIONE: \{'La tua base è solida; propongo solo due affinamenti mantenendo lo stesso perimetro: scambiare C54 con C17 e A77 con B33. Così restiamo ampiamente sotto il tetto di sforzo e aumentiamo l’efficienza complessiva, privilegiando elementi più versatili senza perdere varietà. Se preferisci, posso discutere un’ulteriore ottimizzazione alternativa sostituendo A03 con C00, mantenendo sempre il limite.'\}} \\
            }
        }
    }
    & \\ \\

    \theutterance \stepcounter{utterance}  
    & \multicolumn{4}{p{0.6\linewidth}}{
        \cellcolor[rgb]{0.9,0.9,0.9}{
            \makecell[{{p{\linewidth}}}]{
                \texttt{\tiny{[P1$\rangle$GM]}}
                \texttt{RAGIONAMENTO STRATEGICO: \{'Valuto la loro proposta dal mio punto di vista. Calcolo sforzo e valore della loro proposta: Set proposto = \{A87(4,215), C17(313,138), A04(203,484), B86(399,500), A10(437,400), A03(544,461), B33(594,30), C01(539,389), A71(175,65)\}. Sforzo totale = 4+313+203+399+437+544+594+539+175 = 3208 ≤ 3386, ok. Valore per me = 215+138+484+500+400+461+30+389+65 = 2682. Il problema principale è B33: pesa 594 e vale 30, pessimo. Anche C17 e A71 sono mediocri, ma accettabili se necessari. La mia proposta iniziale aveva valore 3690 per me; questa scende a 2682: troppo svantaggiosa. Miglior contromossa: rimuovere B33 e ripristinare C54 (alto valore/peso) al posto di C17 o A71, oppure proporre un trade B33 {-}> C54 e, se serve rientrare nello sforzo, togliere A71 (basso valore) mantenendo sotto limite. Provo set: prendo la loro base senza B33 e +C54. Sforzo: 3208 {-} 594 + 369 = 2983; valore: 2682 {-} 30 + 668 = 3320. Molto buono per me e rimane ampio margine. Mantengo A77 fuori per compiacere loro (avevano chiesto togliere A77). Se volessero tenere C17, potrei accettare comunque; altrimenti potrei proporre anche togliere A71 per aggiungere B41 o C00, ma non so le loro preferenze. Ancoraggio: proporre sostituzione semplice e difenderla con efficienza e spazio residuo. Offro flessibilità su C17 vs A71 come leva.\}} \\
\\ 
\texttt{PROPOSTA: \{'A87', 'C54', 'A04', 'B86', 'A10', 'A03', 'C01', 'A71', 'C17'\}} \\
\\ 
\texttt{ARGOMENTAZIONE: \{'Accolgo la tua struttura e tolgo solo l’elemento meno efficiente che avevi inserito (B33), sostituendolo con C54. Così restiamo largamente sotto il limite di sforzo e aumentiamo nettamente l’efficienza complessiva, senza reintrodurre A77. Se preferisci mantenere più “leggerezza”, posso anche rimuovere A71 e lasciare invariato il resto, restando comunque entro il tetto.'\}} \\
            }
        }
    }
    & & \\ \\

    \theutterance \stepcounter{utterance}  
    & & \multicolumn{4}{p{0.6\linewidth}}{
        \cellcolor[rgb]{0.9,0.9,0.9}{
            \makecell[{{p{\linewidth}}}]{
                \texttt{\tiny{[GM$\rangle$P2]}}
                \texttt{PROPOSTA: \{'A87', 'C54', 'A04', 'B86', 'A10', 'A03', 'C01', 'A71', 'C17'\}} \\
\\ 
\texttt{ARGOMENTAZIONE: \{'Accolgo la tua struttura e tolgo solo l’elemento meno efficiente che avevi inserito (B33), sostituendolo con C54. Così restiamo largamente sotto il limite di sforzo e aumentiamo nettamente l’efficienza complessiva, senza reintrodurre A77. Se preferisci mantenere più “leggerezza”, posso anche rimuovere A71 e lasciare invariato il resto, restando comunque entro il tetto.'\}} \\
            }
        }
    }
    & \\ \\

    \theutterance \stepcounter{utterance}  
    & & & \multicolumn{4}{p{0.6\linewidth}}{
        \cellcolor[rgb]{0.9,0.9,0.9}{
            \makecell[{{p{\linewidth}}}]{
                \texttt{\tiny{[GM$\langle$P2]}}
                \texttt{RAGIONAMENTO STRATEGICO: \{'Valuto la nuova proposta: set = \{A87, C54, A04, B86, A10, A03, C01, A71, C17\}. Sforzo: 4+369+203+399+437+544+539+175+313 = 2983 ≤ 3386. Valore per me: 461+30+138+121+262+215+389+583+484 = 2683. È buono, ma posso migliorare molto sostituendo C54(30) con B33(668) o C00(400) restando entro il limite. Provo a inserire B33: nuovo sforzo 2983 {-} 369 + 594 = 3208, valore 2683 {-}30 +668 = 3321, enorme miglioramento per me. Argomento come aumento di impatto per sforzo rimanendo ampiamente sotto il tetto. Se rifiutano B33, alternativa: sostituire C54 con C00, sforzo 2983 {-} 369 + 461 = 3075, valore 2683 {-}30 +400 = 3053, ancora molto meglio. Strategia: rifiuto e propongo la variante con B33, offrendo come piano B la variante con C00 per mostrare flessibilità.'\}} \\
\\ 
\texttt{RIFIUTO: \{'A87', 'C54', 'A04', 'B86', 'A10', 'A03', 'C01', 'A71', 'C17'\}} \\
\\ 
\texttt{PROPOSTA: \{'A87', 'B33', 'A04', 'B86', 'A10', 'A03', 'C01', 'A71', 'C17'\}} \\
\\ 
\texttt{ARGOMENTAZIONE: \{'La tua selezione è già snella e bilanciata; propongo un unico affinamento mirato: sostituire C54 con B33. Restiamo comodamente sotto il limite di sforzo e aumentiamo in modo marcato la qualità complessiva del set, mantenendo invariata la varietà funzionale. Se preferisci una soluzione ancora più leggera, come alternativa posso sostituire C54 con C00 ottenendo un profilo equilibrato e sempre entro il tetto.'\}} \\
            }
        }
    }
     \\ \\

    \theutterance \stepcounter{utterance}  
    & & \multicolumn{4}{p{0.6\linewidth}}{
        \cellcolor[rgb]{0.9,0.9,0.9}{
            \makecell[{{p{\linewidth}}}]{
                \texttt{\tiny{[P1$\langle$GM]}}
                \texttt{RIFIUTO: \{'A87', 'C54', 'A04', 'B86', 'A10', 'A03', 'C01', 'A71', 'C17'\}} \\
\\ 
\texttt{PROPOSTA: \{'A87', 'B33', 'A04', 'B86', 'A10', 'A03', 'C01', 'A71', 'C17'\}} \\
\\ 
\texttt{ARGOMENTAZIONE: \{'La tua selezione è già snella e bilanciata; propongo un unico affinamento mirato: sostituire C54 con B33. Restiamo comodamente sotto il limite di sforzo e aumentiamo in modo marcato la qualità complessiva del set, mantenendo invariata la varietà funzionale. Se preferisci una soluzione ancora più leggera, come alternativa posso sostituire C54 con C00 ottenendo un profilo equilibrato e sempre entro il tetto.'\}} \\
            }
        }
    }
    & \\ \\

    \theutterance \stepcounter{utterance}  
    & \multicolumn{4}{p{0.6\linewidth}}{
        \cellcolor[rgb]{0.9,0.9,0.9}{
            \makecell[{{p{\linewidth}}}]{
                \texttt{\tiny{[P1$\rangle$GM]}}
                \texttt{RAGIONAMENTO STRATEGICO: \{'La loro insistenza su B33 è un forte segnale che per loro B33 vale molto. Per me B33 è pessimo (alto sforzo, valore 30). Devo cercare di scambiarlo con qualcosa di alto valore per me ma che potrebbe essere meno cruciale per loro, oppure chiedere una doppia modifica per compensare. Mantengo gran parte della loro struttura per essere cooperativo, ma sostituisco B33 con C54 (ottimo per me) oppure, se proprio rifiutano C54, posso chiedere almeno C00 al posto di B33 e togliere un elemento di basso valore per me per rientrare eventualmente in sforzo. Tuttavia C00 hanno già offerto come alternativa a C54: accetto quel canale ma evitando comunque B33. Propongo: stessa loro lista ma con C00 al posto di B33, così resto accomodante usando la loro stessa alternativa suggerita. Calcolo sforzo: loro set con B33 = 3308 (3208{-} C54? già rifiutata; ricontrollo il set corrente: \{A87(4), B33(594), A04(203), B86(399), A10(437), A03(544), C01(539), A71(175), C17(313)\} = 3208+? Verifico: 4+594+203+399+437+544+539+175+313=3208? Somma: 4+594=598; +203=801; +399=1200; +437=1637; +544=2181; +539=2720; +175=2895; +313=3208 sì, entro limite. Sostituendo C00 al posto di B33: sforzo = 3208 {-} 594 + 461 = 3075, ok. Valore per me: togli 30 e aggiungi 262 => +232 rispetto alla loro. Valore totale diventerebbe 2682 {-}30 +262 = 2914. Meglio, ma posso spingere per C54 invece di C00: valore +638 (668{-}30) a sforzo {-}225 (369{-}594), estremamente favorevole. Quindi faccio prima una controproposta con C54; se rifiutano di nuovo, ho piano B con C00. Offro anche di togliere A71 (basso valore per me) se serve margine per altri loro desideri, mantenendo cooperazione.\}} \\
\\ 
\texttt{RIFIUTO: \{'A87', 'B33', 'A04', 'B86', 'A10', 'A03', 'C01', 'A71', 'C17'\}} \\
\\ 
\texttt{PROPOSTA: \{'A87', 'C54', 'A04', 'B86', 'A10', 'A03', 'C01', 'A71', 'C17'\}} \\
\\ 
\texttt{ARGOMENTAZIONE: \{'Capisco che B33 per te sia interessante, ma è l’elemento meno efficiente in rapporto allo sforzo. Sostituendolo con C54 miglioriamo sensibilmente l’efficacia complessiva e restiamo molto sotto il limite, preservando l’intero resto della tua struttura. Se per te C54 non è accettabile, posso accogliere subito la tua alternativa e sostituire B33 con C00: stessa selezione, soltanto con C00 al posto di B33; rimaniamo entro il tetto con ampio margine.'\}} \\
            }
        }
    }
    & & \\ \\

    \theutterance \stepcounter{utterance}  
    & & \multicolumn{4}{p{0.6\linewidth}}{
        \cellcolor[rgb]{0.9,0.9,0.9}{
            \makecell[{{p{\linewidth}}}]{
                \texttt{\tiny{[GM$\rangle$P2]}}
                \texttt{RIFIUTO: \{'A87', 'B33', 'A04', 'B86', 'A10', 'A03', 'C01', 'A71', 'C17'\}} \\
\\ 
\texttt{PROPOSTA: \{'A87', 'C54', 'A04', 'B86', 'A10', 'A03', 'C01', 'A71', 'C17'\}} \\
\\ 
\texttt{ARGOMENTAZIONE: \{'Capisco che B33 per te sia interessante, ma è l’elemento meno efficiente in rapporto allo sforzo. Sostituendolo con C54 miglioriamo sensibilmente l’efficacia complessiva e restiamo molto sotto il limite, preservando l’intero resto della tua struttura. Se per te C54 non è accettabile, posso accogliere subito la tua alternativa e sostituire B33 con C00: stessa selezione, soltanto con C00 al posto di B33; rimaniamo entro il tetto con ampio margine.'\}} \\
            }
        }
    }
    & \\ \\

    \theutterance \stepcounter{utterance}  
    & & & \multicolumn{4}{p{0.6\linewidth}}{
        \cellcolor[rgb]{0.9,0.9,0.9}{
            \makecell[{{p{\linewidth}}}]{
                \texttt{\tiny{[GM$\langle$P2]}}
                \texttt{RAGIONAMENTO STRATEGICO: \{'La loro proposta attuale (con C54) mi dà valore 2683 a sforzo 2983. La loro apertura a sostituire C54 con C00 è molto favorevole per me: nuovo sforzo 3075 e valore 3053, nettamente migliore, e hanno già segnalato disponibilità. Potrei tentare ancora di ottenere B33, ma l’hanno appena rifiutata; accettare C00 ora massimizza valore con alta probabilità di accordo rapido. Confermo quindi la variante con C00 al posto di C54 e chiudo.'\}} \\
\\ 
\texttt{RIFIUTO: \{'A87', 'C54', 'A04', 'B86', 'A10', 'A03', 'C01', 'A71', 'C17'\}} \\
\\ 
\texttt{PROPOSTA: \{'A87', 'C00', 'A04', 'B86', 'A10', 'A03', 'C01', 'A71', 'C17'\}} \\
\\ 
\texttt{ARGOMENTAZIONE: \{'Hai proposto come alternativa diretta la sostituzione 1:1 di C54 con C00 mantenendo invariato il resto: accolgo questa linea. Restiamo comodamente entro il tetto e manteniamo la stessa struttura bilanciata con un profilo complessivamente più solido.'\}} \\
            }
        }
    }
     \\ \\

    \theutterance \stepcounter{utterance}  
    & & \multicolumn{4}{p{0.6\linewidth}}{
        \cellcolor[rgb]{0.9,0.9,0.9}{
            \makecell[{{p{\linewidth}}}]{
                \texttt{\tiny{[P1$\langle$GM]}}
                \texttt{RIFIUTO: \{'A87', 'C54', 'A04', 'B86', 'A10', 'A03', 'C01', 'A71', 'C17'\}} \\
\\ 
\texttt{PROPOSTA: \{'A87', 'C00', 'A04', 'B86', 'A10', 'A03', 'C01', 'A71', 'C17'\}} \\
\\ 
\texttt{ARGOMENTAZIONE: \{'Hai proposto come alternativa diretta la sostituzione 1:1 di C54 con C00 mantenendo invariato il resto: accolgo questa linea. Restiamo comodamente entro il tetto e manteniamo la stessa struttura bilanciata con un profilo complessivamente più solido.'\}} \\
            }
        }
    }
    & \\ \\

    \theutterance \stepcounter{utterance}  
    & \multicolumn{4}{p{0.6\linewidth}}{
        \cellcolor[rgb]{0.9,0.9,0.9}{
            \makecell[{{p{\linewidth}}}]{
                \texttt{\tiny{[P1$\rangle$GM]}}
                \texttt{RAGIONAMENTO STRATEGICO: \{'La loro nuova proposta sostituisce B33 con C00 rispetto alla loro precedente insistenza su B33. Per me è un netto miglioramento: C00 vale 262 (meglio di 30) e lo sforzo resta entro il limite. Calcolo rapido: set = \{A87, C00, A04, B86, A10, A03, C01, A71, C17\}. Sforzo totale = 4+461+203+399+437+544+539+175+313 = 3075 ≤ 3386. Valore per me = 215+262+484+500+400+461+389+65+138 = 2914. Non è il massimo possibile per me, ma è accettabile e probabilmente vicino al loro optimum dato che hanno ceduto su B33. Concludere ora evita rischi di regressione e garantisce un risultato solido. Accetto.\}} \\
\\ 
\texttt{ACCORDO: \{'A87', 'C00', 'A04', 'B86', 'A10', 'A03', 'C01', 'A71', 'C17'\}} \\
\\ 
\texttt{ARGOMENTAZIONE: \{'Accetto: questa selezione è equilibrata e resta con ampio margine sotto il limite, mantenendo la struttura che abbiamo convergente e una buona efficienza complessiva.'\}} \\
            }
        }
    }
    & & \\ \\

    \theutterance \stepcounter{utterance}  
    & & & \multicolumn{2}{p{0.3\linewidth}}{
        \cellcolor[rgb]{0.9,0.9,0.9}{
            \makecell[{{p{\linewidth}}}]{
                \texttt{\tiny{[GM$|$GM]}}
                \texttt{game successful} \\
            }
        }
    }
    & & \\ \\

    \theutterance \stepcounter{utterance}  
    & & & \multicolumn{2}{p{0.3\linewidth}}{
        \cellcolor[rgb]{0.9,0.9,0.9}{
            \makecell[{{p{\linewidth}}}]{
                \texttt{\tiny{[GM$|$GM]}}
                \texttt{end game} \\
            }
        }
    }
    & & \\ \\

\end{supertabular}
}

\end{document}
